\input{../tex-inputs/latex-leseansicht-vorspann}

\section[Gustav Schwarzkopf an Arthur Schnitzler, 6. 7. 1901]{L04308 Gustav Schwarzkopf an Arthur Schnitzler, 6. 7. 1901}
\nopagebreak\mylabel{L04308v}
\rehead{ }\normalsize\beginnumbering\briefempfaengerindex{Schnitzler, Arthur@\textsc{Schnitzler, Arthur}!zzzSchwarzkopf, Gustav@\emph{von Gustav Schwarzkopf}!1901-07-061@{6. 7. 1901}|(be}
\toendnotes[C]{\smallbreak\pagebreak[2]}
\correspDesc{Versand  durch Gustav Schwarzkopf am 6. 7. 1901 in Wien
\newline{}Erhalt  durch Arthur Schnitzler im Zeitraum [7. 7. 1901 – 11. 7. 1901?] in Wien}\toendnotes[C]{\smallbreak}
\Standort{CUL, Schnitzler, B 96a.}
\physDesc{Brief, 1 Blatt, 4 Seiten, 2954 Zeichen
\newline{}Handschrift: schwarze Tinte, deutsche Kurrent}\toendnotes[C]{\smallbreak}
\pstart
           \raggedleft{}{\pb}6. VII. 01\pend
           \vspace{0.5em}
\pstart
           Lieber Arthur, mein Bruder Max\pwindex{Schwarzkopf, Max 12.\,6.\,1857 Wien – 14.\,4.\,1928 ebd.@\textsc{Schwarzkopf, Max} (12.\,6.\,1857 Wien – 14.\,4.\,1928 ebd.), \emph{Rechtsanwalt}|pw}, der wieder 8 Tage unwol war, iſt geſtern auf 4 Wochen abgereiſt; daraus
               ergibt{ }ſich – Sie ke{\geminationn}en ja{ }ſchon die Einteilung unſeres
               Hauſes – daß ich jetzt nicht gut fort ka{\geminationn}, überdies{ }ſoll
               jetzt entſchieden werden ob unſer Haus\oindex{Wien@\textbf{Wien}!I., Innere Stadt@\textbf{I., Innere Stadt}!Tiefer Graben 23@\textbf{Tiefer Graben 23}, \emph{Wohngebäude}|pwv}{ }ſtehen bleiben und unſere Wohnung hergerichtet werden{ }ſoll oder nicht,
               was wieder – in meiner Eigenſchaft als Hausfrau – meine Anweſenheit erfordert. Sie{ }ſehen alſo –{ }ſo liebenswürdige Formen Ihr »Egoismus« auch anni{\geminationm}t, um mich aufzufordern, und{ }ſo gern ich{ }ſelbſt wieder
               einmal ein paar Tage mit Ihnen verbringen möchte, es geht nicht, wenigſtens nicht im
               Juli, vielleicht im Auguſt, we{\geminationn} Sie doch vielleicht in
               die Nähe des Wörther See\oindex{Wörthersee@\textbf{Wörthersee}, \emph{See}|pw}’s ko{\geminationm}en und ich – den {\pb}»dringenden Einladungen« folgend,
               vielleicht doch einige Tage dort verbringe. – Über die Guſtel\pwindex{Schnitzler, Arthur 15.\,5.\,1862 Wien – 21.\,10.\,1931 ebd.@\textsc{Schnitzler, Arthur} (15.\,5.\,1862 Wien – 21.\,10.\,1931 ebd.), \emph{Schriftsteller, Mediziner}!Lieutenant Gustl. Novelle@\strich\emph{Lieutenant Gustl. Novelle}|pw}-Angelegenheit hab ich nichts geleſen – ich war damals
               grad in der Brühl\oindex{Brühl@\textbf{Brühl}, \emph{Tal}|pw} – als den \label{K_L04308-1v}\edtext{Leitartikel\pwindex{Wien, 20. Juni@\emph{Wien, 20. Juni}|pwv}}{\lemma{\textnormal{\emph{Leitartikel}}}\Cendnote{\textnormal{[Moriz Benedikt\pwindex{Benedikt, Moriz 27.\,5.\,1849 Kvačice – 18.\,3.\,1920 Wien@\textsc{Benedikt, Moriz} (27.\,5.\,1849 Kvačice – 18.\,3.\,1920 Wien), \emph{Journalist, Herausgeber}|pwk}]: \emph{Wien, 20. Juni}\pwindex{Wien, 20. Juni@\emph{Wien, 20. Juni}|pwk}. In: \emph{Neue Freie Presse}\pwindex{Neue Freie Presse@\emph{Neue Freie Presse}|pwk}, Nr. 13.226, 21.\,6.\,1901, Morgenblatt, S. 1–2.}}}\label{K_L04308-1} der »Preſſe\pwindex{Neue Freie Presse@\emph{Neue Freie Presse}|pwv}« und die »\label{K_L04308-2v}\edtext{Fackel\pwindex{Fackel@\emph{Die Fackel}|pw}\pwindex{Kraus, Karl 28.\,4.\,1874 Jičín – 12.\,6.\,1936 Wien@\textsc{Kraus, Karl} (28.\,4.\,1874 Jičín – 12.\,6.\,1936 Wien), \emph{Schriftsteller, Publizist, Schriftsteller}!(Wieder ein Märtyrer.)@\strich\emph{(Wieder ein Märtyrer.)}|pwv}}{\lemma{\textnormal{\emph{Fackel}}}\Cendnote{\textnormal{Karl Kraus\pwindex{Kraus, Karl 28.\,4.\,1874 Jičín – 12.\,6.\,1936 Wien@\textsc{Kraus, Karl} (28.\,4.\,1874 Jičín – 12.\,6.\,1936 Wien), \emph{Schriftsteller, Publizist, Schriftsteller}|pwk}: \emph{(Wieder ein Märtyrer.)}\pwindex{Kraus, Karl 28.\,4.\,1874 Jičín – 12.\,6.\,1936 Wien@\textsc{Kraus, Karl} (28.\,4.\,1874 Jičín – 12.\,6.\,1936 Wien), \emph{Schriftsteller, Publizist, Schriftsteller}!(Wieder ein Märtyrer.)@\strich\emph{(Wieder ein Märtyrer.)}|pwk} In: \emph{Die Fackel}\pwindex{Fackel@\emph{Die Fackel}|pwk}, Jg. 3, H. 80, Mitte Juni 1901, S. 20–24.}}}\label{K_L04308-2}«. Der Leitartikel\pwindex{Wien, 20. Juni@\emph{Wien, 20. Juni}|pwv} war \uline{echte}{ }N. f. Preſſe\pwindex{Neue Freie Presse@\emph{Neue Freie Presse}|pw} du{\geminationm},
               feig und gemein, – der »brave« Schüler, der nur{ }ſich vor Strafe{ }ſchützen und{ }ſich
               beim Herrn Lehrer einſchmeicheln will – und doch, glauben Sie mir, haben Sie in den
               Augen vieler Leute{ }ſehr an Bedeutung gewo{\geminationn}en, weil Sie
               Gegenſtand eines Leitartikels waren und das Quantum Neid, das Ihnen in dieſen Tagen
               zugedacht war, iſt vielleicht Urſache des Drucks, den Sie in falſcher Diagnoſe – wozu
               wären Sie de{\geminationn} Arzt? – Gebirgsaſthma ne{\geminationn}en und der jetzt hoffentlich{ }ſchon verſchwunden {\pb}iſt. Der Druck nämlich, nicht der
               Neid. Der Reſpekt vor dem Blatt\pwindex{Neue Freie Presse@\emph{Neue Freie Presse}|pwv} und
               beſonders vor dem Gewäſch an \uline{der} Stelle, die{ }ſich{ }ſonſt nur mit Gregorig\pwindex{Gregorig, Josef 27.\,4.\,1846 Bisamberg – 2.\,7.\,1909 Maria Enzersdorf@\textsc{Gregorig, Josef} (27.\,4.\,1846 Bisamberg – 2.\,7.\,1909 Maria Enzersdorf), \emph{Politiker}|pw}, Wolf\pwindex{Wolf, Karl Hermann 27.\,1.\,1862 – 11.\,6.\,1941@\textsc{Wolf, Karl Hermann} (27.\,1.\,1862 – 11.\,6.\,1941), \emph{Schriftsteller, Herausgeber, Abgeordneter}|pw} oder den Reden Wilhelm’s\pwindex{Wilhelm II. von Preußen 27.\,1.\,1859 Potsdam – 4.\,6.\,1941 Gemeente Utrechtse Heuvelrug@\textsc{Wilhelm II. von Preußen} (27.\,1.\,1859 Potsdam – 4.\,6.\,1941 Gemeente Utrechtse Heuvelrug), \emph{Kaiser}|pw}
               beſchäftigt, wurzelt doch noch{ }ſehr tief – trotz Kraus\pwindex{Kraus, Karl 28.\,4.\,1874 Jičín – 12.\,6.\,1936 Wien@\textsc{Kraus, Karl} (28.\,4.\,1874 Jičín – 12.\,6.\,1936 Wien), \emph{Schriftsteller, Publizist, Schriftsteller}|pw}. Die N. f. Preſſe\pwindex{Neue Freie Presse@\emph{Neue Freie Presse}|pw} hat alſo gewiß
               Ihren Dank erwartet und eigentlich blieb Ihnen zu nichts andres übrig nach den
               Geſetzen der Höflichkeit. Daß Sie{ }ſich darüber ärgern, es getan zu haben, begreife
               ich{ }ſehr gut. Es iſt{ }ſo wie hie und da bei \textsc{Soupers}. Man hat{ }ſchlecht gegeſſen und{ }ſich{ }ſehr gelangweilt aber beim Fortgehen dankt man doch der
               Hausfrau für den genußreichen Abend und{ }ſchämt{ }ſich dabei, während man es ausſpricht.
               Wenigſtens ich habe mir früher dieſe Regieſpeſen – das Schämen – gemacht. Jetzt denke
               ich nicht mehr. {\pb}Die zwei Nu{\geminationm}ern\pwindex{Unrecht an Aaron Schnitzler@\emph{Das Unrecht an Aaron Schnitzler}|pwv} »Kikeriki\pwindex{Kikeriki@\emph{Kikeriki}|pw}« werde ich heute
               beſorgen. Sehen Sie, wie es{ }ſich verlohnt, gemeint zu{ }ſein: Wären Sie im »Kikeriki\pwindex{Kikeriki@\emph{Kikeriki}|pw}«, geprieſen worden, Sie würden
               vielleicht nicht einmal \uline{eine} Nu{\geminationm}er kaufen laſſen. – Richard\pwindex{Beer-Hofmann, Richard 11.\,7.\,1866 Wien – 26.\,9.\,1945 New York City@\textsc{Beer-Hofmann, Richard} (11.\,7.\,1866 Wien – 26.\,9.\,1945 New York City), \emph{Schriftsteller}|pw} und Frau\pwindex{Beer-Hofmann, Paula 25.\,2.\,1879 Wien – 30.\,10.\,1939 Zürich@\textsc{Beer-Hofmann, Paula} (25.\,2.\,1879 Wien – 30.\,10.\,1939 Zürich)|pwv}
               habe ich hier in Wien\oindex{Wien@\textbf{Wien}, \emph{Verwaltungsgebiet}|pw} geſprochen; die Frau\pwindex{Beer, Rosa 20.\,7.\,1847 Borek Fałęcki – 1.\,7.\,1901 Wien@\textsc{Beer, Rosa} (20.\,7.\,1847 Borek Fałęcki – 1.\,7.\,1901 Wien)|pwv} seines Vaters \textsc{Beer\pwindex{Beer, Hermann 10.\,8.\,1835 Radiměř – 3.\,10.\,1902 Wien@\textsc{Beer, Hermann} (10.\,8.\,1835 Radiměř – 3.\,10.\,1902 Wien), \emph{Rechtsanwalt}|pw}}, die{ }ſchon lange leidend war, iſt plötzlich geſtorben; den Vater\pwindex{Beer, Hermann 10.\,8.\,1835 Radiměř – 3.\,10.\,1902 Wien@\textsc{Beer, Hermann} (10.\,8.\,1835 Radiměř – 3.\,10.\,1902 Wien), \emph{Rechtsanwalt}|pwv} hat er jetzt mit{ }ſich nach Pörtſchach\oindex{Pörtschach am Wörthersee@\textbf{Pörtschach am Wörthersee}|pw} geno{\geminationm}en.–\pend
           
\pstart
           – – Von mir kann ich Ihnen nichts berichten, weder äußere noch innere,
               Erlebniſſe. Im »deutschen Theater\orgindex{Deutsches Theater Berlin@Deutsches Theater Berlin|pw}« war ich nicht,
               dafür habe ich auch »Schall und Rauch\orgindex{Schall und Rauch@Schall und Rauch|pw}« nicht mit
               meinem Beſuch beehrt. Geſchäfte{ }ſollen nicht gut geweſen{ }ſein. –\textcolor{gray}{–}
               Bitte,{ }ſagen Sie den beiden jungen Damen\pwindex{Schnitzler, Olga 17.\,1.\,1882 Wien – 13.\,1.\,1970 Lugano@\textsc{Schnitzler, Olga} (17.\,1.\,1882 Wien – 13.\,1.\,1970 Lugano), \emph{Schauspielerin, Sängerin}|pwv}\pwindex{Steinrück, Elisabeth 19.\,11.\,1885 – 7.\,4.\,1920 Partenkirchen@\textsc{Steinrück, Elisabeth} (19.\,11.\,1885 – 7.\,4.\,1920 Partenkirchen)|pwv} mein beſten Dank für ihre freundlichen Worte und beſtellen Sie meine
               ergebenſten Grüße.\pend
           
\pstart
           Herzlichſt{\\[\baselineskip]}Ihr{\\[\baselineskip]}\spacefill\mbox{Gustav}\pend
           \leftskip=0em{}\selectlanguage{ngerman}\endnumbering\briefempfaengerindex{Schnitzler, Arthur@\textsc{Schnitzler, Arthur}!zzzSchwarzkopf, Gustav@\emph{von Gustav Schwarzkopf}!1901-07-061@{6. 7. 1901}|)be}\mylabel{L04308h}
\begin{anhang}
\end{anhang}\newcommand{\dateiname}{L04308}\newcommand{\titel}{Gustav Schwarzkopf an Arthur Schnitzler, 6. 7. 1901}\newcommand{\editorInnen}{Herausgegeben von Jahnke, SelmaMüller, Martin Anton}\input{../tex-inputs/latex-leseansicht-abspann}
